Bei hohen Temperaturen haben alle festen Elemente nach der Regel von Dulong-Petit dieselbe molare Wärme $C = 3R$ (FoTaBe S.~171). Die spezifische Wärme ist die molare Wärme geteilt durch die molare Masse, $c = C/M = 3R/M$. Damit ist
\begin{align}
 \frac{c_1}{c_2} &= \frac{3R/M_1}{3R/M_2} = \rcF = \frac{\MbO}{\MaO} = \rc
\end{align}
Die spezifischen Wärmen verhalten sich umgekehrt wie die Molmassen: $c_1 : c_2 = \result{\MbO : \MaO}$.
